% =====================================================================================
% La escalera de la estrategia de entrenamiento — figura propia en TikZ.
%
% IDEA: cada peldaño es un BLOQUE que se apoya en el anterior y sobresale hacia la
% derecha. Los bloques se tocan (no hay hueco entre ellos), que es lo que hace que el
% conjunto se lea como una escalera y no como una lista inclinada. Se sube de izquierda
% a derecha porque cada peldaño AÑADE una capacidad sobre el de abajo, y la pregunta en
% gris es justamente lo que ese peldaño aisla frente al anterior.
%
% CODIGO DE COLOR (informativo, no decorativo):
%   azulunir  -> los tres peldaños de REFERENCIA: no son modelos de la comparativa,
%                estan para fijar el suelo contra el que se mide.
%   acentofig -> los TRES MODELOS que se comparan de verdad (§ del capitulo).
%
% ⚠️ ANCHOS MEDIDOS. Caja de texto de la memoria: 15,02 cm.
%    bloque 6,8 cm de texto + 2x0,176 de margen = 7,15 · desplazamiento 1,25 por peldaño
%    -> borde derecho en 5x1,25 + 7,15 = 13,40 · mas la flecha y su rotulo a la izquierda
%    (~0,95) da 14,35 cm. NO ensanchar el bloque sin reducir el desplazamiento.
%
% ⚠️ LA SEGUNDA LINEA DE CADA BLOQUE NO DEBE PARTIRSE. A \scriptsize en 6,8 cm entran
%    ~49 caracteres; las seis estan por debajo. Si se reescribe alguna y pasa de ahi, el
%    bloque crece, deja de medir 1,15 cm y la escalera se descuadra.
%
% ⚠️ SIN TITULO dentro de la imagen: lo pone el \caption de LaTeX (norma del TFE, §1.3).
%
% ⚠️ \shorthandoff{<>} es obligatorio: `babel` en español hace activos < y >, y eso rompe
%    la sintaxis de puntas de flecha de TikZ.
% =====================================================================================
\begingroup
\shorthandoff{<>}%
\hyphenpenalty=10000\exhyphenpenalty=10000\relax

\begin{tikzpicture}[
    >={Stealth[length=2mm]},
    font=\footnotesize,
    peldano/.style = {draw=#1, line width=0.9pt, fill=#1!8, rounded corners=1pt,
                      text width=6.8cm, minimum height=1.15cm, align=left,
                      inner sep=5pt, anchor=south west},
    nota/.style    = {font=\scriptsize, color=gray},
  ]

  % --- los seis peldaños ---------------------------------------------------------------
  % Referencias: fijan el suelo, no entran en la comparativa.
  \node[peldano=azulunir] at (0.00, 0.00)
    {\textbf{1. No mitigar}\\[1pt]
     {\scriptsize\color{gray} el suelo: si no se bate, el método no sirve}};

  \node[peldano=azulunir] at (1.25, 1.15)
    {\textbf{2. Factor de supervivencia constante}\\[1pt]
     {\scriptsize\color{gray} la media de entrenamiento · ¿aporta saberla?}};

  \node[peldano=azulunir] at (2.50, 2.30)
    {\textbf{3. Ridge con una sola variable}\\[1pt]
     {\scriptsize\color{gray} la duración total · ¿aporta saber cuánto dura?}};

  % Los tres que se comparan.
  \node[peldano=acentofig] at (3.75, 3.45)
    {\textbf{4. Ridge con todas las variables}\\[1pt]
     {\scriptsize\color{gray} ¿aporta solo la linealidad?}};

  \node[peldano=acentofig] at (5.00, 4.60)
    {\textbf{5. Random Forest}\\[1pt]
     {\scriptsize\color{gray} ¿aporta la no linealidad?}};

  \node[peldano=acentofig] at (6.25, 5.75)
    {\textbf{6. Red de grafos}\\[1pt]
     {\scriptsize\color{gray} ¿aporta la estructura del circuito?}};

  % --- el eje que dice hacia donde se sube ----------------------------------------------
  \draw[->, azulunir, line width=1pt] (-0.55, 0) -- (-0.55, 6.90)
        node[midway, rotate=90, above=1pt, font=\scriptsize, text=black]
        {capacidad del modelo};

  % --- leyenda, en el hueco que deja la escalera abajo a la derecha ----------------------
  % Va aqui y no arriba a la izquierda por dos razones: arriba quedaba pegada al \caption,
  % y abajo a la derecha rellena el unico hueco grande que deja la escalera.
  \draw[draw=azulunir, line width=0.9pt, fill=azulunir!8]  (7.60, 0.72) rectangle (7.88, 1.00);
  \node[nota, anchor=west] at (7.98, 0.86) {peldaños de referencia: fijan el suelo};

  \draw[draw=acentofig, line width=0.9pt, fill=acentofig!8] (7.60, 0.22) rectangle (7.88, 0.50);
  \node[nota, anchor=west] at (7.98, 0.36) {los tres modelos que se comparan};

\end{tikzpicture}%
\endgroup
